% compendium/energy-commodities/commodity-curves.tex
% Entry migrated on 2026-10-16 from the legacy macro chapter (slice 37, 02_transmission_maps.tex), owned by
% commodity-storage-curves. Gliner (2014) is cited by chapter, section and PDF page of the ebook.
% Dispositions in notes/migration-slices.md, slice 37.

\section{Curves and inventories}

\begin{qfdefinition}[Contango, backwardation and ending stocks: storage and financing costs, scarcity premia, seasonal crop curves, the stocks-to-use ratio, and roll yield stated carefully]\label{def:contango-backwardation-stocks}
\qffield{Statement} Gliner: commodities are driven by supply (producers, costs, yield, weather, politics) and demand (population, development, alternative energy) \citep[chap.~10, secs.\ ``Supply Drivers'' and ``Demand Drivers'', PDF pp.~293--296]{s7-gliner}. \emph{Ending stock} ``effectively measures the inventory of a particular commodity''; crops are harvested once or twice a year while demand is continuous, and the stock carried over becomes the next season's beginning balance (sec.\ ``Ending Stock'', p.~296, Figure~10.1). \emph{Contango} is near-term futures trading below later ones, from storage, transport and the opportunity cost of capital (gold in a vault); \emph{backwardation} is near-term futures at a premium, as in a supply shortage (cocoa during an Ivory Coast export stoppage) or just before a harvest when stocks are depleted; a crop curve can show both, contango after harvest and backwardation before the next (sec.\ ``Contango and Backwardation'', pp.~297--298, Figure~10.2, corn as of 30 September 2013).

\qffield{Notation} Balance sheet: beginning stocks $+$ production $+$ imports $-$ consumption $-$ exports $=$ ending stocks; stocks-to-use $=$ ending stocks $/$ total use. Full carry: $F=S e^{(r+u-y)T}$ with storage cost $u$ and convenience yield $y$ (\hyperref[res:cost-of-carry-forward-price]{cost-of-carry forward price}).

\qffield{Conditions} The legacy identity ``ending stocks $=$ beginning stocks $+$ production $-$ consumption'' holds for the world, not for a country, which also has trade; with imports of 5 and exports of 12 on a balance of 10, 100 and 95, ending stocks are 8, not 15 (\nolinkurl{notes/artefacts/multi-asset-systematic-macro/slice37-checks.py}, section~16). Low stocks relative to use support prices and make them sensitive to small supply changes, which is the storage model the pathway estimates; the legacy's ``most important driver of commodity prices'' is an overstatement of it.

\qffield{Practice} The monthly WASDE report \citep[catalogue entry only, not read]{physical-usda-wasde-2026} gives USDA's U.S.\ and world supply-and-demand balances. ``Normal state for storable commodities like gold'' (legacy contango) fits monetary metals with large above-ground stocks; it is not a general rule for storable commodities, which can be backwardated when inventories are scarce.

\qffield{Pitfalls} ``Contango creates negative roll yield; backwardation creates positive roll yield'' (legacy) is true of the roll component holding the curve fixed, not of total returns, and it says nothing on its own about a risk premium: whether backwardation reflects scarcity or compensation for risk is the pathway's question. \emph{Own bounded test} (batch of 2026-10-17; not a reading, and not evidence beyond its sample). On EIA's public NYMEX series (contracts 1, 2 and~4, which EIA stopped updating on 5 April 2024) and EIA inventories, with backwardation $B_t=\ln(C_{1,t}/C_{2,t})$ at month-ends, inventories as the log deviation from the same week's mean of the previous five years (lagged five days), and the next month's return of holding contract~2, the hypotheses and pass rules were fixed before computing: WTI, August 1987 to March 2024 (439 months), backwardation rises as inventories fall ($\beta=-0.112$, Newey--West $t=-5.9$; with $C_1/C_4$, $t=-6.5$), and Henry Hub, 2015 to March 2024 (111 months), likewise ($\beta=-0.124$, $t=-4.4$): the scarcity link is supported in both. The pre-specified premium test is not supported: the slope of next-month returns on $B_t$ is $+0.05$ ($t=0.1$) for WTI and $-0.19$ ($t=-0.6$) for gas, inventories add nothing, and WTI's out-of-sample $R^2$ against the historical mean is $-2.2$ percent (Clark--West $t=-0.4$; gas too short for the pre-set 120-month start). Exploratory only, added after the first run: WTI's mean monthly return was $+2.2$ percent in backwardated months against $-0.3$ in contango (difference $t=3.4$), but the linear and out-of-sample forms fail, two months exceed $\pm50$ percent, and gas shows the opposite sign ($-5.0$ against $-0.4$). Costs of two ticks per roll barely move these figures. Script, protocol and outputs: \nolinkurl{notes/artefacts/commodity-storage-curves/backwardation_test.py}, \nolinkurl{backwardation_results.txt} and the derived monthly panel \nolinkurl{backwardation_monthly.csv} in the same folder. Limits: two U.S. energy contracts; monthly horizon; no other commodities; returns from EIA's continuous contract series under the expiry mapping stated in the script.

\qffield{See also} \hyperref[res:cost-of-carry-forward-price]{cost-of-carry forward price}; \hyperref[def:futures-hedge-basis]{futures hedges and basis}; \hyperref[def:commodity-indices-risk-return]{commodity indices}; \hyperref[def:agricultural-markets]{agricultural markets}.

\qffield{Further reading} The competitive storage model and convenience yield, the pathway's readings.

\qffield{Index} contango; backwardation; ending stocks; beginning stocks; stocks-to-use ratio; supply and demand balance; WASDE; storage costs; convenience yield; seasonal curve; crop year; roll yield

\qffield{Status} Verification: partly checked 2026-10-16 (Gliner's text, chap.~10, PDF pp.~291--298; the trade-adjusted balance computed in the slice's check script; batch of 2026-10-17: own test on EIA data, reported in Pitfalls). Stability: draft. Currency: timeless concepts; Gliner's examples as of 2013; test data to March 2024.
\end{qfdefinition}
